\chapter{Empatyzacja}

\section{Wprowadzenie do fazy empatyzacji}

Empatyzacja stanowi pierwszy i kluczowy etap procesu projektowego opartego na metodyce
Design Thinking. Jej celem jest dogłębne zrozumienie potrzeb, oczekiwań, obaw oraz
motywacji wszystkich osób zaangażowanych lub dotkniętych planowanym wdrożeniem systemu
informatycznego. W przypadku projektu systemu EMR dla Szpitala im. Bajgla we Wrocławiu
szczególnie istotne jest uwzględnienie szerokiego spektrum interesariuszy od dyrekcji
placówki, przez personel medyczny i techniczny, aż po pacjentów oraz potencjalnych
użytkowników zewnętrznych.

Zidentyfikowano cztery
główne grupy interesariuszy: kluczowych graczy, użytkowników (z podziałem na głównych i potencjalnych) oraz interesariuszy uśpionych oraz obserwatorów. Dla każdej z grup
opracowano charakterystykę persony wraz z wynikającymi z niej implikacjami projektowymi.

\section{Mapa interesariuszy}

Na potrzeby projektu interesariuszy pogrupowano według ich roli i wpływu na wdrożenie systemu:

\begin{itemize}
    \item \textbf{Kluczowi gracze} - osoby o wysokim wpływie na projekt i wysokim poziomie zaangażowania: dyrektor szpitala, zarząd szpitala, pracownik IT
    \item \textbf{Główni użytkownicy} - osoby bezpośrednio korzystające z systemu w codziennej pracy tzn. pracownicy szpitala m.in : lekarz internista, studentka medycyny, pielęgniarka oddziałowa.
    \item \textbf{Potencjalni użytkownicy} - osoby korzystające z systemu okazjonalnie lub w ograniczonym zakresie - pacjenci. Wyodrębniono również złośliwego użytkownika, przed którym system musi być chroniony.
    \item \textbf{Uśpieni} - interesariusze o potencjalnie dużym wpływie, lecz aktualnie niezaangażowani: CEO startupu technologicznego, media, ministerstwo zdrowia.
    \item \textbf{Obserwatorzy} - grupa interesariuszy zewnętrznych mających minimalny wpływ na projekt: administracja publiczna, audytorzy, fundacje, ksiądz, policja
\end{itemize}

\section{Analiza person}

\subsection{Dyrektor szpitala}

\textbf{Profil:} Mężczyzna, 63 lata. Doświadczony administrator z wieloletnim stażem w ochronie zdrowia, żonaty, aktywny dziadek.

\textbf{Cechy charakteru:} Miły, ciepły, troskliwy, nowoczesny i optymistyczny. Otwarty na zmiany, o ile służą pacjentom i pracownikom.

\textbf{Nastawienie do projektu:} Pozytywne. Dyrektor jest inicjatorem i głównym sponsorem projektu, dostrzega potrzebę cyfryzacji i chce poprawić standardy działania placówki.

\textbf{Oczekiwania wobec systemu:}
\begin{itemize}
    \item Zgodność z obowiązującymi standardami prawnymi i regulacjami ochrony danych,
    \item Realne usprawnienie codziennej pracy lekarzy i personelu,
    \item Intuicyjność dla wszystkich grup wiekowych użytkowników,
    \item Wysoki poziom bezpieczeństwa danych medycznych.
\end{itemize}

\textbf{Obawy:}
\begin{itemize}
    \item Możliwość pojawienia się sztucznych problemów technicznych,
    \item sprzeciw części lekarzy wobec wdrożenia,
    \item ryzyko wycieku danych pacjentów,
\end{itemize}

\textbf{Implikacje projektowe:} System musi być możliwy do przedstawienia jako realna korzyść dla personelu, a nie dodatkowe obciążenie. Kluczowe jest przygotowanie rzetelnej prezentacji wartości projektu dla interesariuszy sceptycznych. Dyrektor powinien zostać włączony w kampanię komunikacyjną poprzedzającą wdrożenie.

\subsection{Zarząd szpitala}

\textbf{Profil:} Zarząd szpitala jest organem kolegialnym składającym się z kilku osób pełniących funkcje kierownicze najwyższego szczebla. W typowej strukturze dużego szpitala publicznego w Polsce w jego skład wchodzą: zastępca dyrektora ds. medycznych (lekarz z wieloletnim stażem klinicznym i administracyjnym), zastępca dyrektora ds. finansowych (ekonomista lub prawnik z doświadczeniem w finansach publicznych), główna księgowa oraz kierownicy kluczowych pionów organizacyjnych. Zarząd działa w ścisłej współpracy z dyrektorem, jednak jako ciało kolegialne podejmuje decyzje strategiczne, w tym o zatwierdzeniu budżetu projektu.
Z perspektywy profilu demograficznego są to osoby w wieku 48–58 lat, z długim stażem w sektorze publicznej ochrony zdrowia. Większość z nich budowała karierę w czasach, gdy cyfryzacja szpitali była jeszcze pieśnią przyszłości, co kształtuje ich ostrożne, analityczne podejście do inwestycji technologicznych.

\textbf{Nastawienie do projektu:} Warunkowo pozytywne. Członkowie zarządu rozumieją nieuchronność cyfryzacji. Jednak znają przykłady wdrożenia nowych systemów IT w innych placówkach: opóźnione, przekraczające budżet, wywołujące opór personelu. To doświadczenie rodzi instynktowną ostrożność.

\textbf{Oczekiwania wobec systemu:}
\begin{itemize}
    \item bezpieczeństwo prawne - system musi być w pełni zgodny z RODO, ustawą o prawach pacjenta i Rzeczniku Praw Pacjenta,
    \item ciągłość operacyjna - szpital działa 24 godziny na dobę, 7 dni w tygodniu i nie może pozwolić sobie na przestój systemu krytycznego.
    \item kontrola kosztów - zarząd oczekuje szczegółowego kosztorysu, jasnego określenia kosztów utrzymania w perspektywie co najmniej 5 lat oraz zdefiniowania warunków ewentualnego odstąpienia od umowy z dostawcą.
    \item wymierne efekty - zarząd będzie monitorował czy system rzeczywiście skraca czas pracy administracyjnej personelu.
\end{itemize}

\textbf{Obawy:}
\begin{itemize}
    \item ryzyko przekroczenia budżetu i harmonogramu,
    \item opór personelu,
    \item ewentualne roszczenia pacjentów w przypadku błędu systemowego mającego wpływ na proces leczenia.
\end{itemize}

\textbf{Implikacje projektowe:} zarząd będzie oceniał projekt głównie przez pryzmat kosztów, zgodności prawnej oraz wpływu na efektywność operacyjną placówki. Kluczowe będzie przygotowanie szczegółowego kosztorysu, harmonogramu oraz analizy zwrotu z inwestycji. Zarząd będzie wymagał formalnej dokumentacji zgodności z regulacjami RODO oraz dyrektywami UE dotyczącymi eHealth.

\subsection{Pracownik IT}

\textbf{Profil:} Mężczyzna, 48 lat. Zna dobrze dyrektora, co czyni go kluczowym pośrednikiem. Życie prywatne złożone: żonaty, jednocześnie rozwiedziony, syn z poprzedniego związku. Motto: \emph{,,Nie da się - nie ruszaj''}.

\textbf{Cechy charakteru:} Sceptyczny, ostrożny, pedantyczny, realistyczny. Ceni stabilność ponad innowacyjność.

\textbf{Nastawienie do projektu:} Negatywne. Spodziewa się znacznego wzrostu nakładu pracy, dodatkowych zgłoszeń od personelu i odpowiedzialności za system, który sam nie projektował.

\textbf{Oczekiwania wobec systemu:}
\begin{itemize}
    \item Pełna stabilność i niezawodność,
    \item Jasna i kompletna dokumentacja techniczna,
    \item Realne wsparcie techniczne od dostawcy.
\end{itemize}


\textbf{Obawy:} 
\begin{itemize}
    \item Lawina pytań od personelu,
    \item brak szkoleń,
    \item niedopracowanie systemu w momencie wdrożenia,
    \item błędy trudne do diagnozowania.
\end{itemize}

\textbf{Implikacje projektowe:} Pracownik IT musi zostać aktywnie włączony w projekt już na etapie projektowania architektury. Niezbędne jest przekazanie mu szczegółowej dokumentacji technicznej, jasnych procedur eskalacji problemów oraz gwarancji wsparcia ze strony dostawcy. Jego akceptacja jest warunkiem powodzenia wdrożenia ze względu na bliską relację z dyrekcją.

\subsection{Starszy lekarz internista}

\textbf{Profil:} Mężczyzna, 60 lat. Doświadczony specjalista z ugruntowaną pozycją zawodową. Żona księgowa, dwóch dorosłych synów. Czyta prasę drukowaną, ma skłonność do krytycznej oceny nowych trendów. Motto: \emph{,,Wszystko musi się zgadzać''}.

\textbf{Cechy charakteru:} Uparty, sceptyczny wobec technologii, przekonany o własnej racji, oparty na wieloletnich doświadczeniach klinicznych.

\textbf{Nastawienie do projektu:} Negatywne. Uważa projekt za utrudnienie pracy, a nie jej usprawnienie.

\textbf{Oczekiwania wobec systemu:}
\begin{itemize}
    \item System ma działać szybciej niż tradycyjny zapis odręczny,
    \item Dane mają uzupełniać się automatycznie wszędzie, gdzie to możliwe,
    \item Możliwość wydrukowania każdego dokumentu w dowolnym momencie.
\end{itemize}

\textbf{Obawy:} 
\begin{itemize}
    \item Utrata danych w razie awarii systemu
    \item dostęp nieuprawnionych osób do jego pacjentów
    \item wyparcie z roli przez młodszych lekarzy sprawniejszych technologicznie,
    \item konieczność długiego klikania zamiast leczenia.
\end{itemize}

\textbf{Implikacje projektowe:} Interfejs musi być maksymalnie uproszczony, minimalnie trzy kliknięcia do dowolnej informacji. Konieczne jest zapewnienie podglądu i wydruku każdego dokumentu. System nie może spowalniać lekarza podczas obchodu. Zalecane jest przeprowadzenie indywidualnych szkoleń dla tej grupy wiekowej przed wdrożeniem.

\subsection{Studentka medycyny}

\textbf{Profil:} Kobieta, 25 lat. Mieszka z partnerem w wynajmowanym mieszkaniu.

\textbf{Cechy charakteru:} Wytrwała, ambitna, uporządkowana, empatyczna. Sprawnie porusza się w środowiskach cyfrowych.

\textbf{Nastawienie do projektu:} Pozytywne. Widzi w systemie szansę na naukę i lepsze zrozumienie procesów klinicznych.

\textbf{Oczekiwania wobec systemu:}
\begin{itemize}
    \item Intuicyjny i nowoczesny interfejs,
    \item Możliwość dostępu do danych pacjenta i obserwacji decyzji lekarza prowadzącego,
    \item Szeroki zakres dostępnych informacji.
\end{itemize}

\textbf{Obawy:} System zbyt skomplikowany jak na jej poziom uprawnień, konieczność proszenia personelu o pomoc przy każdej operacji.

\textbf{Implikacje projektowe:} Należy zaprojektować dedykowany poziom dostępu dla stażystów i studentów, umożliwiający obserwację bez możliwości modyfikacji danych. Interfejs powinien zawierać kontekstowe podpowiedzi i opisy, wspierające uczenie się systemu.

\subsection{Pielęgniarka oddziałowa}

\textbf{Profil:} Kobieta, 38 lat. Mężatka, dwoje dzieci w wieku szkolnym. Aktywna fizycznie, biegaczka i rowerzystka. Na co dzień korzysta ze smartfona i mediów społecznościowych. Motto: \emph{,,Liczy się pacjent, nie papierologia''}.

\textbf{Cechy charakteru:} Zorganizowana, konkretna, zorientowana na efektywność. Bywa niecierpliwa wobec zbędnych procedur. Liderka nieformalna, jej nastawienie wpływa na postawy innych pielęgniarek.

\textbf{Nastawienie do projektu:} Ambiwalentne. Rozumie potrzebę cyfryzacji, lecz obawia się równoległego działania dwóch systemów w okresie przejściowym. Jeśli system okaże się intuicyjny, stanie się jego ambasadorką.

\textbf{Oczekiwania wobec systemu:}
\begin{itemize}
    \item Szybkie wprowadzanie danych przy minimalnej liczbie kliknięć,
    \item Czytelny widok zleceń i zadań dla każdego pacjenta,
    \item Powiadomienia push o pilnych zleceniach,
    \item Działanie w trybie offline lub przy słabym sygnale Wi-Fi,
    \item Integracja z dotychczasową dokumentacją.
\end{itemize}

\textbf{Obawy:} 
\begin{itemize}
    \item Czas poświęcony na wpisywanie danych kosztem opieki nad pacjentem
    \item Awarie podczas dyżuru, 
    \item Szkolenie przeprowadzone na ostatnią chwilę, 
    \item Odpowiedzialność za błędy systemowe.
\end{itemize}
\textbf{Implikacje projektowe:} System musi posiadać tryb offline z synchronizacją po odzyskaniu połączenia. Projektując interfejs mobilny, należy kłaść nacisk na ergonomię użytkowania tabletu. Powiadomienia o zleceniach są wymaganiem krytycznym. Pielęgniarkę warto zaangażować jako testera w fazie pilotażowej, ponieważ jej pozytywna ocena ułatwi wdrożenie w całym zespole pielęgniarskim.

\subsection{Starszy pacjent}

\textbf{Profil:} Mężczyzna, 72 lata. Wdowiec, dwoje dorosłych dzieci w innym mieście. Czas spędza na działce, ogląda telewizję, rozmawia z sąsiadami.

\textbf{Cechy charakteru:} Spokojny, cierpliwy, ale łatwo się frustruje gdy coś nie działa intuicyjnie. Nieufny wobec nowych technologii. Potrzebuje jasnych, prostych instrukcji.

\textbf{Nastawienie do projektu:} Sceptyczne - nierozumie po co to wszystko, skoro wszystko działało bez komputerów. 

\textbf{Oczekiwania wobec systemu:}
\begin{itemize}
    \item Bardzo prosty interfejs z dużą czcionką i wysokim kontrastem,
    \item Możliwość przeglądania wyników badań bez wychodzenia z domu,
    \item Powiadomienia SMS o nowych wynikach,
    \item Możliwość wydruku wyników.
\end{itemize}

\textbf{Obawy:} 
\begin{itemize}
    \item Brak zdolności samodzielnej obsługi systemu,
    \item Obawa przed wyciekiem danych osobowych i medycznych,
    \item Niezrozumiałość wyników bez komentarza lekarza, 
    \item Brak wsparcia technicznego na miejscu.
\end{itemize}
Brak zdolności samodzielnej obsługi systemu, obawa przed wyciekiem danych osobowych i medycznych, niezrozumiałość wyników bez komentarza lekarza, brak wsparcia technicznego na miejscu.

\textbf{Implikacje projektowe:} Portal pacjenta musi spełniać wymagania dostępności cyfrowej tj. duże przyciski, czytelna typografia (min. 16 pt), wysoki kontrast, prosty język. Powiadomienia SMS muszą być domyślnie włączone. Przy wynikach badań powinna pojawiać się interpretacja w języku niespecjalistycznym lub odsyłacz do lekarza.
\subsection{Pacjent w średnim wieku}

\textbf{Profil:} Kobieta, 35 lat. Zamężna, dwoje dzieci w wieku przedszkolnym. Pracuje jako
specjalistka ds. marketingu w trybie hybrydowym. Na co dzień aktywnie korzysta ze smartfona,
mediów społecznościowych oraz aplikacji bankowych i zakupowych. Motto:
\emph{,,Chcę to załatwić szybko i wracać do swoich spraw.''}

\textbf{Cechy charakteru:} Pragmatyczna, zorientowana na efektywność, niecierpliwa wobec
rozwiązań nieintuicyjnych. Przyzwyczajona do wysokiego standardu UX aplikacji komercyjnych takich jak aplikacji bankowej lub e-commerce. Samodzielna i rzadko
prosi o pomoc techniczną, o ile interfejs działa zgodnie z oczekiwaniami.

\textbf{Nastawienie do projektu:} Pozytywne i praktyczne. Nie interesuje jej technologia sama
w sobie, lecz konkretna oszczędność czasu. Ważne jest także brak konieczności brania urlopu na odbiór
wyników, unikanie kolejek i możliwość sprawdzenia dokumentacji między spotkaniami w pracy.

\textbf{Oczekiwania wobec systemu:}
\begin{itemize}
    \item Szybki i wygodny dostęp do wyników badań przez aplikację mobilną lub stronę WWW,
    \item Powiadomienia push lub SMS w momencie pojawienia się nowych wyników,
    \item Podgląd historii leczenia i zleceń lekarskich w jednym miejscu,
    \item Możliwość pobrania dokumentów w formacie PDF na potrzeby drugiej opinii lub ubezpieczyciela,
    \item Przejrzysty, nowoczesny interfejs działający równie sprawnie na telefonie co na komputerze.
\end{itemize}

\textbf{Obawy:} 
\begin{itemize}
    \item Wyciek danych medycznych lub ich wykorzystanie przez firmy ubezpieczeniowe i pracodawców,
    \item Brak responsywności interfejsu na urządzeniach mobilnych,
    \item Konieczność rejestrowania się przez wiele kroków i zapamiętywania kolejnego hasła,
    \item Ryzyko, że wyniki będą widoczne w systemie z opóźnieniem i nadal trzeba będzie dzwonić do szpitala.
\end{itemize}

\textbf{Implikacje projektowe:} Portal pacjenta musi być w pełni responsywny i zoptymalizowany
pod kątem urządzeń mobilnych. Rejestracja i logowanie powinny umożliwiać uwierzytelnianie
przez profil zaufany lub aplikację mObywatel, co eliminuje potrzebę zapamiętywania
dodatkowych danych dostępowych. Wyniki badań muszą być publikowane w systemie
automatycznie, niezwłocznie po ich zatwierdzeniu przez laboratorium. Dokumenty powinny
być dostępne do pobrania w formacie PDF. Polityka prywatności musi być sformułowana
w zrozumiałym, niespecjalistycznym języku, a cel przetwarzania każdej kategorii danych jasno opisany przy pierwszym logowaniu.

\subsection{Złośliwy użytkownik - przedstawiciel usług pogrzebowych}

\textbf{Profil:} Mężczyzna, 45 lat. Rozwiedziony, dwoje dzieci. Często spędza czas za kierownicą, towarzyski, niekiedy spożywa alkohol. Prowadzi działalność gospodarczą. Motto: \emph{,,Kto pierwszy, ten lepszy''}.

\textbf{Cechy charakteru:} Otwarty, ambitny, sprytny, przedsiębiorczy. Potrafi nawiązywać relacje z personelem.

\textbf{Nastawienie do projektu:} Optymistyczne - liczy na ułatwienie pozyskiwania klientów dzięki słabym zabezpieczeniom systemu.

\textbf{Oczekiwania (nieuprawnione):} 
\begin{itemize}
    \item Łatwy dostęp do danych pacjentów, 
    \item Możliwość podglądu dokumentacji przez przekupienie personelu lub wykorzystanie luk w systemie.
\end{itemize}

\textbf{Obawy:} Silne uwierzytelnianie, zaawansowane mechanizmy kontroli dostępu, audyt logowań.

\textbf{Implikacje projektowe:} Dane medyczne należą do kategorii danych szczególnie wrażliwych w rozumieniu RODO. System musi wdrażać: uwierzytelnianie dwuskładnikowe (2FA), pełny rejestr zdarzeń i logowań, automatyczne wylogowanie po czasie bezczynności, szyfrowanie danych zarówno w spoczynku, jak i podczas transmisji. Personel medyczny powinien zostać przeszkolony z zakresu ochrony danych i zasad nieudzielania dostępu osobom nieuprawnionym.

\subsection{CEO małego startupu technologicznego}

\textbf{Profil:} Kobieta lub mężczyzna, 32 lata. Zapracowany, rodzice nie wiedzą, czym dokładnie się zajmuje. Motto: \emph{,,Pójdzie jak z płatka''}.

\textbf{Cechy charakteru:} Optymistka, nieskończona energia, tendencja do ignorowania szczegółów, wiara w skalowalność każdego rozwiązania.

\textbf{Rola w projekcie:} Potencjalny inwestor i dostawca technologii.

\textbf{Nastawienie do projektu:} Euforyczne - widzi szansę na zmianę rynku ochrony zdrowia.

\textbf{Oczekiwania:} Szybkie wdrożenie, zbieranie jak największej ilości danych.

\textbf{Obawy:} Biurokracja, długie procedury, wymogi prawne, RODO.

\textbf{Implikacje projektowe:} Należy jasno określić zakres przetwarzanych danych i podstawy prawne ich zbierania. Minimalizacja danych (privacy by design) jest wymogiem RODO - system nie może gromadzić danych w zakresie szerszym niż niezbędny do realizacji celu medycznego. W relacjach z dostawcami zewnętrznymi konieczne są umowy powierzenia przetwarzania danych (art. 28 RODO).

\subsection{Media}

\textbf{Profil:} Interesariusz uśpiony o potencjalnie dużym wpływie na odbiór projektu przez opinię publiczną.

\textbf{Implikacje projektowe:} Media mogą odegrać ważną rolę w przypadku incydentu bezpieczeństwa lub protestu personelu. Rekomenduje się przygotowanie strategii komunikacji zewnętrznej na etapie wdrożenia oraz wyznaczenie rzecznika projektu.

\subsection{Ministerstwo Zdrowia}

\textbf{Profil:} Interesariusz regulacyjny o wysokim potencjalnym wpływie.

\textbf{Implikacje projektowe:} System musi być zgodny z wymogami P1 (Platformy e-Zdrowia) oraz polskim prawem o dokumentacji medycznej (Ustawa o prawach pacjenta, Rozporządzenie MZ w sprawie rodzajów i zakresu dokumentacji medycznej). Wdrożenie powinno zostać skonsultowane z Centrum e-Zdrowia.

\section{Synteza wniosków - zdefiniowane wymagania funkcjonalne}

Na podstawie przeprowadzonej analizy person zidentyfikowano następujące kluczowe wymagania funkcjonalne i niefunkcjonalne, które muszą być uwzględnione w projekcie systemu:

\paragraph{Wymagania wynikające z analizy użytkowników:}
\begin{itemize}
    \item \textbf{Dostęp rolowy} - każda grupa użytkowników (lekarz, pielęgniarka, laborant, stażysta, pacjent, administrator) otrzymuje oddzielny zestaw uprawnień dostosowany do zakresu obowiązków.
    \item \textbf{Mobilność} - aplikacja musi działać na tabletach i smartfonach; personel oddziałowy nie ma stałego dostępu do komputerów stacjonarnych.
    \item \textbf{Tryb offline} - ze względu na słabą infrastrukturę sieciową na niektórych oddziałach system musi obsługiwać pracę bez połączenia z serwerem i synchronizować dane po odzyskaniu dostępu do sieci.
    \item \textbf{Szybkość operacji} - dowolna informacja o pacjencie powinna być dostępna w nie więcej niż 3 krokach; krytyczne jest to dla lekarza podczas obchodu oraz pielęgniarki podczas dyżuru.
    \item \textbf{Portal pacjenta} - prosty, dostępny cyfrowo (WCAG 2.1 AA) interfejs z możliwością przeglądania wyników, powiadomieniami SMS oraz opcją wydruku.
    \item \textbf{Powiadomienia i alerty} - automatyczne powiadamianie personelu o nowych wynikach badań i pilnych zleceniach.
    \item \textbf{Możliwość wydruku} - każdy dokument musi być dostępny do wydruku w standardowym formacie.
\end{itemize}

\paragraph{Wymagania wynikające z analizy bezpieczeństwa:}
\begin{itemize}
    \item \textbf{Uwierzytelnianie dwuskładnikowe (2FA)} - obowiązkowe dla wszystkich kont personelu medycznego,
    \item \textbf{Pełny rejestr zdarzeń (audit log)} - każde logowanie, wylogowanie i zmiana danych musi być zapisywana z identyfikatorem użytkownika i znacznikiem czasu,
    \item \textbf{Szyfrowanie danych} - zarówno w spoczynku (AES-256), jak i podczas transmisji (TLS 1.3),
    \item \textbf{Automatyczne wylogowanie} - po upływie 5 minut bezczynności,
    \item \textbf{Minimalizacja danych} - zbieranie wyłącznie danych niezbędnych do realizacji celu medycznego, zgodnie z art. 5 RODO.
\end{itemize}

\paragraph{Wymagania wynikające z analizy wdrożeniowej:}
\begin{itemize}
    \item \textbf{Szkolenia zróżnicowane grupowo} - odrębne szkolenia dla lekarzy, pielęgniarek, laborantów, administracji i studentów,
    \item \textbf{Dokumentacja techniczna} - kompletna i aktualna dokumentacja dla pracownika IT,
    \item \textbf{Wsparcie powdrożeniowe} - dedykowany helpdesk przez co najmniej 3 miesiące od startu produkcyjnego,
    \item \textbf{Wdrożenie etapowe} - pilotaż na jednym oddziale przed rozszerzeniem na całą placówkę,
    \item \textbf{Zgodność z regulacjami} - RODO, wymogi Centrum e-Zdrowia, Ustawa o prawach pacjenta.
\end{itemize}
